% =============================================================
%  Tugas Individu 03 — Representasi Gambar
%  IF25-40304 · Pembelajaran Mesin Multimodal · ITERA
% =============================================================
\documentclass[12pt, a4paper]{article}

\usepackage{tugas-style}

\renewcommand{\assignmentnumber}{03}
\renewcommand{\doctitle}{Tugas Individu 03}
\renewcommand{\duedate}{\color{ctpBlue}\faCalendarCheck\ \ Tenggat: sesuai jadwal kelas \faCalendarCheck}

\begin{document}

% ── HEADER ──────────────────────────────────────────────────
\makeassignmentheader

\hrule width0.3\textwidth
\revisioninfo{%
  \vhEntry{0.1}{16 September 2026}{I.W.W.}{Draf inisial tugas.}
}
\hrule
\vskip .3in

\tableofcontents
\clearpage

% =============================================================
\section{Tujuan Pembelajaran}\label{sec:objectives}

Setelah menyelesaikan tugas ini, mahasiswa mampu:
\begin{objectives}
  \item Merepresentasikan gambar sebagai tensor RGB berdimensi $H \times W \times C$ beserta metadatanya.
  \item Menggunakan ResNet50 pra-latih sebagai ekstraktor fitur visual (feature extractor).
  \item Menghitung kemiripan kosinus antar-embedding gambar dan menginterpretasikan hasilnya.
  \item Memproyeksikan embedding gambar berdimensi tinggi ke 2D menggunakan PCA.
  \item Menganalisis secara kritis keterbatasan representasi visual dan cosine similarity.
\end{objectives}

% =============================================================
\section{Catatan dan Batasan}\label{sec:notes}
\vspace{-.1cm}

\begin{minipage}{\linewidth}
\begin{bclogo}[couleur=ctpBase, arrondi=0.1, logo=\bccrayon, ombre=false, couleurBord=ctpSurface, epBord=0.6]{Perhatikan hal berikut:}
  \begin{coloredPen}
    \item Tugas ini bersifat \textbf{individu}; tidak diperkenankan bekerja sama atau menyalin jawaban.
    \item Seluruh eksperimen memakai \textbf{inferensi} pada ResNet50 pra-latih --- jangan melakukan \code{fit()}, \code{backward()}, atau pelatihan ulang classifier. Pada Soal 2 Anda \textbf{menghitung konvolusi secara manual} dan memverifikasinya dengan \code{torch.nn.functional.conv2d}; keduanya operasi maju (\emph{forward}), bukan pelatihan.
    \item Fokus tugas adalah mengamati dan menganalisis representasi visual, bukan melatih model baru.
    \item Siapkan gambar lokal yang dapat dibandingkan secara visual; unduh weights ResNet50 lebih awal.
  \end{coloredPen}
\end{bclogo}
\end{minipage}

% =============================================================
\section{Perangkat Lunak dan Peralatan yang Diperlukan}\label{sec:requiredsw}

{\emergencystretch=2em
\begin{itemize}[itemsep=2pt,parsep=0pt,topsep=2pt,partopsep=2pt]
  \item[\color{ctpBlue}\faPython] \textbf{Python} 3.8 atau lebih baru.
  \item[\color{ctpBlue}\faFileCode] \textbf{Paket wajib:} \code{torch}, \code{torchvision}, \code{Pillow}, \code{matplotlib}, \code{scikit-learn}, \code{numpy}.
  \item[\color{ctpBlue}\faImage] \textbf{Gambar:} dua atau lebih gambar dengan isi yang dapat dibandingkan secara visual.
  \item[\color{ctpBlue}\faDownload] \textbf{Internet:} saat pertama kali mengunduh weights ResNet50.
\end{itemize}}

\begin{codewindow}{setup.sh}
\begin{lstlisting}[style=pycode]
pip install torch torchvision pillow matplotlib scikit-learn
\end{lstlisting}
\end{codewindow}

\begin{minipage}{\linewidth}
\begin{bclogo}[couleur=ctpMantle, arrondi=0.1, logo=\bcinfo, ombre=false, couleurBord=ctpSurface, epBord=0.6]{Batasan Cakupan}
  Seluruh tugas menggunakan inferensi pada model ResNet50 pra-latih. Jangan melakukan
  \code{fit()}, \code{backward()}, atau pelatihan ulang classifier. Pada Soal 2 Anda
  \textbf{menghitung konvolusi secara manual} dan memverifikasinya dengan
  \code{torch.nn.functional.conv2d}; keduanya operasi maju (\emph{forward}), bukan
  pelatihan. Fokus tugas adalah mengamati dan menganalisis representasi visual.
\end{bclogo}
\end{minipage}

% =============================================================
\section{Pernyataan Masalah}\label{sec:intro}

Tugas ini melengkapi materi Pertemuan 03 tentang representasi gambar, CNN, ResNet, dan transfer learning. Anda akan melakukan eksperimen inferensi menggunakan model ResNet50 pra-latih untuk memverifikasi empat gagasan utama:
\begin{enumerate}
  \item Gambar memiliki struktur spasial dan tiga kanal warna RGB, bukan sekadar daftar angka independen.
  \item Konvolusi dan pooling menyaring pola lokal (tepi, tekstur) dengan jumlah parameter yang jauh lebih sedikit daripada lapisan \emph{fully connected}.
  \item CNN membangun representasi hierarkis dari pola lokal menuju fitur visual yang lebih abstrak.
  \item Backbone pra-latih dapat digunakan sebagai ekstraktor fitur visual tanpa melatih ulang model.
\end{enumerate}

Rincian dan persyaratan pengerjaan diberikan pada bagian berikutnya.

% =============================================================
\section{Persyaratan}\label{sec:requirements}

Perhatikan setiap persyaratan pada soal-soal berikut.

% ------------------------------------------------------------
\subsection{Soal 1 --- Gambar sebagai Tensor RGB}
\label{sec:soal1}

Gambar RGB dapat direpresentasikan sebagai tensor dengan bentuk $H \times W \times C$, dengan $H$ sebagai tinggi, $W$ sebagai lebar, dan $C=3$ sebagai kanal merah, hijau, dan biru.

\subsubsection{Instruksi}
\label{sec:soal1-instruksi}

\textbf{(a)} Pilih dua gambar lokal dengan ukuran atau rasio aspek yang berbeda. Muat gambar menggunakan Pillow, lalu tampilkan:
\begin{itemize}
  \item ukuran asli setiap gambar;
  \item mode warna;
  \item jumlah kanal;
  \item rasio aspek.
\end{itemize}

\textbf{(b)} Tampilkan kedua gambar menggunakan \code{matplotlib} dan jelaskan perbedaan visualnya secara singkat.

\textbf{(c)} Jawab dalam 3--5 kalimat:
\begin{itemize}
  \item Mengapa gambar tidak cukup direpresentasikan sebagai daftar angka tanpa struktur spasial?
  \item Apa konsekuensi perbedaan ukuran dan rasio aspek sebelum gambar diproses CNN?
  \item Mengapa gambar RGB memiliki $C=3$?
\end{itemize}

\subsubsection{Contoh Kerangka Kode}
\label{sec:soal1-kode}

\begin{codewindow}{soal1\_tensor.py}
\begin{lstlisting}[style=pycode]
from pathlib import Path
from PIL import Image
import matplotlib.pyplot as plt

paths  = [Path("data/image_a.jpg"), Path("data/image_b.jpg")]
images = [Image.open(path).convert("RGB") for path in paths]

for path, image in zip(paths, images):
    width, height = image.size
    print(path.name, image.size, image.mode)
    print("channels:", len(image.getbands()))
    print("aspect ratio:", round(width / height, 3))

fig, axes = plt.subplots(1, 2, figsize=(10, 4))
for axis, path, image in zip(axes, paths, images):
    axis.imshow(image)
    axis.set_title(path.name)
    axis.axis("off")
plt.tight_layout()
plt.show()
\end{lstlisting}
\end{codewindow}

% ------------------------------------------------------------

% ------------------------------------------------------------
\clearpage
\subsection{Soal 2 --- Konvolusi 2D dan Pooling}
\label{sec:soal2}

Slide 8 menuliskan operasi konvolusi sebagai $y(i,j) = \sum_{u,v} x(i+u, j+v)\, k(u,v)$, sedangkan slide 9 menjelaskan \emph{stride} dan \emph{padding}. Keduanya sering dipahami sebagai rumus hafalan. Soal ini menuntutnya dihitung pada satu contoh kecil sehingga terlihat bahwa setiap nilai \emph{feature map} hanyalah penjumlahan hasil kali piksel lokal dengan kernel.

Perhatikan tambalan (\emph{patch}) abu-abu $6 \times 6$ berikut. Baris atas bernilai 0 (gelap) dan baris bawah bernilai 1 (terang), sehingga terdapat satu \textbf{tepi horizontal} di antara baris ketiga dan keempat.

\begin{center}
\ttfamily
\begin{tabular}{@{}l@{}}
patch = \\
{[}0 0 0 0 0 0{]} \\
{[}0 0 0 0 0 0{]} \\
{[}0 0 0 0 0 0{]} \\
{[}1 1 1 1 1 1{]} \\
{[}1 1 1 1 1 1{]} \\
{[}1 1 1 1 1 1{]} \\
\end{tabular}
\end{center}

Gunakan dua kernel Sobel $3 \times 3$:

\begin{equation*}
k_y = \begin{bmatrix} -1 & -2 & -1 \\ 0 & 0 & 0 \\ 1 & 2 & 1 \end{bmatrix},
\qquad
k_x = \begin{bmatrix} -1 & 0 & 1 \\ -2 & 0 & 2 \\ -1 & 0 & 1 \end{bmatrix}
\end{equation*}

\subsubsection{Instruksi}
\label{sec:soal2-instruksi}

\textbf{(a)} Hitung \textbf{satu} nilai keluaran $y(1,0)$ secara manual: ambil jendela $3 \times 3$ dari \code{patch[1:4, 0:3]}, kalikan elemen demi elemen dengan $k_y$, lalu jumlahkan. Tuliskan langkah aritmetikanya, bukan hanya hasil akhirnya.

\textbf{(b)} Hitung seluruh \emph{feature map} untuk $k_y$ dengan perulangan bersarang (\code{for} ganda) memakai \code{numpy}, lalu bandingkan dengan hasil \code{torch.nn.functional.conv2d}. Tampilkan kedua matriks dan selisih maksimumnya.

\textbf{(c)} Untuk masukan $6 \times 6$ dan kernel $3 \times 3$, hitung ukuran keluaran dengan rumus
\begin{equation*}
H_{out} = \left\lfloor \frac{H - K + 2P}{S} \right\rfloor + 1
\end{equation*}
untuk tiga konfigurasi: $(P{=}0, S{=}1)$, $(P{=}1, S{=}1)$, dan $(P{=}1, S{=}2)$. Verifikasi setiap nilai dengan memeriksa \code{.shape} hasil \code{conv2d}.

\textbf{(d)} Terapkan max pooling $2 \times 2$ pada \emph{feature map} $k_y$. Laporkan bentuk sebelum dan sesudahnya, serta nilai yang bertahan.

\textbf{(e)} Jalankan langkah (b) dengan $k_x$ sebagai ganti $k_y$. Jelaskan mengapa hasilnya nol di seluruh posisi, padahal pada (b) tidak demikian. Kaitkan jawaban Anda dengan gagasan bahwa kernel \textbf{memilih} pola tertentu.

\textbf{(f)} Bandingkan jumlah parameter sebuah \emph{fully connected layer} yang memetakan masukan $64 \times 64 \times 3$ ke 64 unit dengan sebuah lapisan konvolusi yang memakai 64 filter $3 \times 3$ (abaikan bias). Jelaskan mengapa \emph{parameter sharing} membuat konvolusi lebih layak dipakai pada gambar.

\subsubsection{Contoh Kerangka Kode}
\label{sec:soal2-kode}

\begin{codewindow}{soal2\_konvolusi.py}
\begin{lstlisting}[style=pycode]
import numpy as np
import torch
import torch.nn.functional as F

patch = np.array([[0, 0, 0, 0, 0, 0],
                  [0, 0, 0, 0, 0, 0],
                  [0, 0, 0, 0, 0, 0],
                  [1, 1, 1, 1, 1, 1],
                  [1, 1, 1, 1, 1, 1],
                  [1, 1, 1, 1, 1, 1]], dtype=np.float32)

kernel_y = np.array([[-1, -2, -1],
                     [ 0,  0,  0],
                     [ 1,  2,  1]], dtype=np.float32)

# (a) satu nilai secara manual
window = patch[1:4, 0:3]
print(window)
print("y(1,0) =", (window * kernel_y).sum())

# (b) seluruh feature map + verifikasi
H, W = patch.shape
K = kernel_y.shape[0]
feature = np.zeros((H - K + 1, W - K + 1), dtype=np.float32)

for i in range(feature.shape[0]):
    for j in range(feature.shape[1]):
        feature[i, j] = (patch[i:i+K, j:j+K] * kernel_y).sum()

tensor = torch.from_numpy(patch)[None, None]
weight = torch.from_numpy(kernel_y)[None, None]
library = F.conv2d(tensor, weight)[0, 0].numpy()

print(feature)
print(library)
print("selisih maks:", np.abs(feature - library).max())
\end{lstlisting}
\end{codewindow}

\subsubsection{Pemeriksaan Mandiri}
\label{sec:soal2-hasil}

Tiga nilai berikut hanya untuk memeriksa apakah kode Anda berjalan benar. \textbf{Ukuran keluaran pada (c) dan penjelasan pada (e) sengaja tidak dicantumkan} karena keduanya termasuk yang dinilai.

\begin{center}
\small
\begin{tabular}{@{}L{4.2cm} L{8.5cm}@{}}
  \toprule \textbf{Pemeriksaan} & \textbf{Nilai yang benar} \\
  \midrule
  Hasil (a) & $y(1,0) = 4$ \\
  \emph{Feature map} $k_y$ & bernilai $4$ pada dua baris tengah; nol pada baris teratas dan terbawah \\
  Hasil (d) & $4 \times 4$ menjadi $2 \times 2$, seluruh nilainya $4$ \\
  Selisih manual vs \code{conv2d} & $0.0$ \\
  \bottomrule
\end{tabular}
\end{center}

Jika selisih maksimum pada (b) tidak nol, periksa urutan indeks jendela (\code{patch[i:i+K, j:j+K]}) dan pastikan kernel tidak dibalik secara tidak sengaja.

\subsection{Soal 3 --- ResNet50 sebagai Feature Extractor}
\label{sec:soal2-resnet}

ResNet50 yang telah dilatih pada ImageNet dapat digunakan sebagai backbone untuk menghasilkan embedding visual. Classifier terakhir dilepas sehingga keluaran model berupa vektor fitur berdimensi 2048.

\subsubsection{Instruksi}
\label{sec:soal2resnet-instruksi}

\textbf{(a)} Muat weights ResNet50 dan ubah classifier terakhir menjadi \code{torch.nn.Identity()}.

\textbf{(b)} Gunakan transformasi resmi dari weights model. Jalankan inferensi pada dua gambar yang dipilih pada Soal 1.

\textbf{(c)} Tampilkan bentuk keluaran setiap gambar. Pastikan hasilnya berupa vektor \code{(2048,)} setelah dimensi batch dihilangkan.

\begin{codewindow}{soal2\_resnet50.py}
\begin{lstlisting}[style=pycode]
import torch
from torchvision.models import ResNet50_Weights, resnet50

weights  = ResNet50_Weights.DEFAULT
backbone = resnet50(weights=weights)
backbone.fc = torch.nn.Identity()
backbone.eval()

preprocess = weights.transforms()

features = []
with torch.inference_mode():
    for image in images:
        tensor  = preprocess(image).unsqueeze(0)
        feature = backbone(tensor).squeeze(0)
        features.append(feature)
        print("feature shape:", tuple(feature.shape))
\end{lstlisting}
\end{codewindow}

\textbf{(d)} Jawab dalam 3--5 kalimat:
\begin{itemize}
  \item Mengapa classifier terakhir dilepas?
  \item Mengapa mode \code{eval()} dan \code{torch.inference\_mode()} sesuai untuk tugas ini?
  \item Mengapa keluaran backbone berupa vektor, bukan lagi matriks piksel?
\end{itemize}

% ------------------------------------------------------------
\subsection{Soal 4 --- Kemiripan Kosinus Antar-Gambar}
\label{sec:soal3-kosinus}

Dua gambar dapat dibandingkan melalui arah vektor embedding. Kemiripan kosinus dihitung dengan:
\[
\operatorname{sim}(\mathbf{a}, \mathbf{b}) =
\frac{\mathbf{a} \cdot \mathbf{b}}
{\lVert\mathbf{a}\rVert\,\lVert\mathbf{b}\rVert}
\]

\subsubsection{Instruksi}
\label{sec:soal3-instruksi}

\textbf{(a)} Hitung kemiripan kosinus antara dua embedding.

\begin{codewindow}{soal3\_cosine.py}
\begin{lstlisting}[style=pycode]
import torch.nn.functional as F

similarity = F.cosine_similarity(
    features[0].unsqueeze(0),
    features[1].unsqueeze(0),
).item()

print(f"cosine similarity: {similarity:.4f}")
\end{lstlisting}
\end{codewindow}

\textbf{(b)} Ulangi eksperimen dengan pasangan gambar kedua. Pilih satu pasangan yang secara visual mirip dan satu pasangan yang berbeda.

\textbf{(c)} Sajikan hasil dalam tabel:

\begin{center}\small
\begin{tabular}{@{} L{2cm} L{4.5cm} C{2.2cm} L{3.5cm} @{}}
  \toprule \textbf{Pasangan} & \textbf{Deskripsi visual} & \textbf{Cosine similarity} & \textbf{Interpretasi} \\
  \midrule
  A & $\ldots$ & $\ldots$ & $\ldots$ \\
  B & $\ldots$ & $\ldots$ & $\ldots$ \\
  \bottomrule
\end{tabular}
\end{center}

\textbf{(d)} Analisis dalam 3--5 kalimat:
\begin{itemize}
  \item Apakah pasangan yang lebih mirip secara visual selalu memiliki nilai cosine lebih tinggi?
  \item Informasi visual apa yang kemungkinan besar memengaruhi nilai tersebut?
  \item Mengapa nilai cosine bukan probabilitas bahwa dua gambar berasal dari kelas yang sama?
\end{itemize}

% ------------------------------------------------------------
\clearpage
\subsection{Soal 5 --- Visualisasi Embedding Gambar}
\label{sec:soal4-visual}

Embedding 2048 dimensi sulit diamati secara langsung. Gunakan PCA untuk memproyeksikan beberapa embedding ke ruang dua dimensi.

\subsubsection{Instruksi}
\label{sec:soal4-instruksi}

\textbf{(a)} Siapkan minimal \textbf{6 gambar} yang terbagi ke dalam \textbf{2 atau 3 kelompok visual}. Contoh kelompok: hewan, kendaraan, makanan, atau pemandangan.

\textbf{(b)} Ekstrak embedding setiap gambar menggunakan pipeline pada Soal 3.

\textbf{(c)} Reduksi embedding ke dua dimensi menggunakan PCA dan buat scatter plot berlabel.

\begin{codewindow}{soal4\_visualisasi.py}
\begin{lstlisting}[style=pycode]
import numpy as np
import matplotlib.pyplot as plt
from sklearn.decomposition import PCA

# features: list tensor berdimensi 2048
matrix      = torch.stack(features).numpy()
coordinates = PCA(n_components=2).fit_transform(matrix)

plt.figure(figsize=(8, 6))
for i, (x, y) in enumerate(coordinates):
    plt.scatter(x, y, s=100)
    plt.annotate(labels[i], (x, y), xytext=(7, 4),
                 textcoords="offset points", fontweight="bold")
plt.title("Proyeksi 2D Embedding Gambar")
plt.xlabel("Komponen Utama 1")
plt.ylabel("Komponen Utama 2")
plt.grid(alpha=0.3)
plt.tight_layout()
plt.savefig("visualisasi_embedding_gambar.png", dpi=150)
plt.show()
\end{lstlisting}
\end{codewindow}

\textbf{(d)} Analisis dalam 3--5 kalimat:
\begin{itemize}
  \item Apakah gambar dalam kelompok yang sama cenderung berdekatan?
  \item Gambar mana yang paling jauh dari kelompoknya?
  \item Mengapa proyeksi PCA dua dimensi tidak selalu mempertahankan seluruh informasi embedding asli?
\end{itemize}

% ------------------------------------------------------------
\subsection{Soal 6 --- Analisis Kritis Representasi Visual}
\label{sec:soal5-analisis}

Jawab pertanyaan berikut dalam bentuk paragraf singkat (masing-masing 5--8 kalimat):

\textbf{(a) Backbone dan Tugas Baru}

ResNet50 dilatih untuk klasifikasi ImageNet, tetapi pada tugas ini digunakan sebagai ekstraktor fitur. Jelaskan mengapa fitur dari model pra-latih dapat digunakan untuk tugas yang berbeda, dan kapan pendekatan ini mungkin kurang sesuai.

\textbf{(b) Struktur Lokal dan Abstraksi}

Hubungkan proses \code{convolution\ $\to$\ pooling\ $\to$\ feature map\ $\to$\ embedding} dengan hierarki representasi dari piksel, tepi, tekstur, hingga bagian objek.

\textbf{(c) Keterbatasan Cosine Similarity}

Jelaskan mengapa dua gambar dapat memiliki nilai cosine yang tinggi meskipun manusia menilai keduanya tidak identik. Pertimbangkan pengaruh dataset ImageNet, preprocessing, latar belakang, dan bias backbone.

% =============================================================
\clearpage
\section{Kriteria Penilaian}\label{sec:evaluation}

Tugas Anda akan dinilai berdasarkan kriteria berikut:

\renewcommand{\arraystretch}{1.5}
\begin{center}\small
\begin{tabular}{|p{11.5cm}|c|}
  \hline
  \thead{\color{ctpBlue} Kriteria} & \thead{\color{ctpBlue}Bobot} \\
  \hline
  Kebenaran teknis --- kode berjalan tanpa error dan pipeline inferensi digunakan dengan tepat. & 35\% \\
  \hline
  Kedalaman analisis --- penjelasan menghubungkan output dengan konsep CNN dan representasi visual. & 30\% \\
  \hline
  Eksperimen --- pemilihan pasangan/kelompok gambar dan interpretasi hasil. & 20\% \\
  \hline
  Kualitas presentasi --- laporan rapi, terstruktur, dan visualisasi informatif. & 15\% \\
  \hline
  \textbf{Total} & \textbf{100\%} \\
  \hline
\end{tabular}
\end{center}

% =============================================================
\section{Apa yang Dikumpulkan}\label{sec:submit}

Seluruh artefak tugas dikemas ke dalam \textbf{satu berkas arsip} dengan format penamaan:
\begin{center}
  \code{mgg03\_tugas\_nim.zip}
\end{center}

dengan \code{nim} diganti oleh Nomor Induk Mahasiswa (NIM) masing-masing tanpa spasi.

\textbf{Artefak yang dapat disertakan di dalam arsip} (dapat lebih dari satu jenis berkas, disesuaikan dengan bentuk penyelesaian Anda):
\begin{borderedsquare}
  \item \code{mgg03\_tugas\_NIM.ipynb} --- apabila penyelesaian dikerjakan menggunakan Jupyter Notebook (termasuk Google Colab). Notebook \textbf{wajib telah dijalankan} sehingga seluruh \emph{output} sel tersimpan di dalam berkas.
  \item \code{mgg03\_tugas\_NIM.py} --- apabila penyelesaian dikerjakan dalam bentuk skrip Python.
  \item \code{mgg03\_tugas\_NIM.pdf} --- apabila analisis ditulis dalam dokumen teks terpisah.
  \item \code{mgg03\_tugas\_xx.jpg} --- apabila terdapat plot atau gambar pendukung tambahan; \code{xx} merupakan nomor urut visualisasi (misalnya \code{01}, \code{02}, dan seterusnya).
\end{borderedsquare}

Setiap berkas menggunakan awalan \code{mgg03\_tugas\_NIM} yang sama, dengan \code{NIM} ditulis secara konsisten (huruf besar/kecil mengikuti NIM resmi). Pastikan seluruh berkas berada tepat di dalam satu arsip \code{.zip} tanpa struktur folder bertingkat yang tidak perlu.

\begin{callout}[ctpBlue]
\textbf{\color{ctpBlue}Catatan Pengerjaan}
\begin{itemize}
  \item \textbf{Soal 1--5:} Wajib dikerjakan oleh semua mahasiswa.
  \item \textbf{Soal 6:} Wajib dikerjakan sebagai analisis konseptual tanpa pelatihan model.
  \item \textbf{Estimasi waktu:} 2--3 jam.
\end{itemize}
\end{callout}

\section*{Referensi}
{\small\color{ctpSubtext}
\begin{enumerate}[label={[\arabic*]}, leftmargin=2em]
  \item He, K., Zhang, X., Ren, S., \& Sun, J. (2016). \emph{Deep Residual Learning for Image Recognition}. CVPR. \url{https://doi.org/10.1109/CVPR.2016.90}
  \item PyTorch Contributors. \emph{ResNet-50 --- Torchvision documentation}. \url{https://pytorch.org/vision/stable/models/generated/torchvision.models.resnet50.html}
  \item Goodfellow, I., Bengio, Y., \& Courville, A. (2016). \emph{Deep Learning}, Bab 9: Convolutional Networks. MIT Press.
\end{enumerate}}

\end{document}
